\chapter{Time-delay systems and the limits of Padé approximations}\label{ch:delay}

This chapter studies the scalar characteristic equation
\begin{equation}
  F(s,T)=D(s)+N(s)\,e^{-sT}=0,\qquad T\ge0,
  \label{eq:qp}
\end{equation}
with real polynomials $D$ of degree $m$ and $N$ of degree $p$. It shows
that the stability of such systems can be decided for every delay without
approximating the delay, and then examines precisely what a Padé
approximation of $e^{-sT}$ preserves and what it loses.

\section{Formulation}

The simplest example is the delay-differential equation
\begin{equation}
  \dot x(t)=-a\,x(t)-b\,x(t-T).
  \label{eq:dde}
\end{equation}
Its solution from $t=0$ is determined only by a whole \emph{history}
$x(\theta)=\varphi(\theta)$, $\theta\in[-T,0]$: on $[0,T]$ the delayed term is
known and \eqref{eq:dde} is an ordinary differential equation; its solution
determines the delayed term on $[T,2T]$, and so on (the \emph{method of
steps} \citep{bellman1963}). The state is a function, and the system is
infinite-dimensional. Exponential solutions $e^{st}$ exist when
$s+a+be^{-sT}=0$. More generally, the loop of Figure~\ref{fig:loop}, a
rational plant $N/D$ with delayed unit feedback, has the characteristic
equation~\eqref{eq:qp}.

\begin{figure}[H]
\centering
\begin{tikzpicture}[
  blk/.style={draw=ink,thick,minimum width=22mm,minimum height=10mm,align=center,fill=soft},
  sum/.style={draw=ink,thick,circle,minimum size=6mm,inner sep=0},
  arr/.style={-{Stealth[length=2.2mm]},thick}]
  \node (in) at (0,0) {$r$};
  \node[sum,right=12mm of in] (s) {};
  \node[blk,right=14mm of s] (G) {$\dfrac{N(s)}{D(s)}$};
  \coordinate[right=16mm of G] (out);
  \node[right=4mm of out] (y) {$y$};
  \node[blk,below=10mm of G,fill=softred,draw=rust] (d) {$e^{-sT}$};
  \draw[arr] (in) -- (s);
  \draw[arr] (s) -- node[above]{$e$} (G);
  \draw[arr] (G) -- (out) -- (y);
  \draw[arr] (out) |- (d);
  \draw[arr] (d) -| node[pos=0.95,right]{$-$} (s);
\end{tikzpicture}
\caption{Loop with a delay in the feedback path. The loop transfer function
is $L(s)=\frac{N(s)}{D(s)}e^{-sT}$; $1+L=0$ is equivalent to~\eqref{eq:qp}.
The delay has unit gain at every frequency and unbounded phase $-\omega T$.}
\label{fig:loop}
\end{figure}

\begin{definition}
The system is \emph{retarded} if $p<m$ and \emph{neutral} if $p=m$.
\end{definition}
Throughout this chapter $p<m$. Three cases illustrate the typical
behaviors (Table~\ref{tab:cases}).

\begin{table}[H]\centering\small
\caption{Case studies; all are stable at $T=0$.}\label{tab:cases}
\begin{tabular}{llll}\toprule
Case & Equation & Model & Behavior as $T$ grows\\\midrule
P1 & $s+1+2e^{-sT}=0$ & $\dot x=-x(t)-2x(t-T)$ & loses stability at $T\approx1.209$\\
P2 & $s+2+e^{-sT}=0$ & $\dot x=-2x(t)-x(t-T)$ & stable for every $T\ge0$\\
P3 & $s^2+0.8s+4+3e^{-sT}=0$ & $\ddot x+0.8\dot x+4x=-3x(t-T)$ & stability switches\\
\bottomrule\end{tabular}
\end{table}

\section{Structure of the spectrum}

$F(\cdot,T)$ is entire and, for $T>0$, has infinitely many roots. For
retarded systems they form chains that move to the left: balancing
$|D(s)|\sim|s|^m$ with $|N(s)e^{-sT}|\sim|s|^pe^{-T\Real s}$ gives
\[
  \Real s\approx-\frac{m-p}{T}\ln|s|,\qquad |s|\to\infty .
\]
Hence every half-plane $\Real s>c$ contains finitely many roots, the
\emph{spectral abscissa} $\alpha(T)=\max\{\Real s:F(s,T)=0\}$ is well
defined, and the zero solution of a linear retarded system is exponentially
stable if and only if $\alpha(T)<0$ \citep{hale1993,michiels2014}.
Figure~\ref{fig:polemaps} shows the spectra of the three cases together
with the poles of Padé models (Section~\ref{sec:pade}).

\begin{figure}[H]
\centering
\includegraphics[width=\textwidth]{pole_maps.pdf}
\caption{Roots of~\eqref{eq:qp} ($\times$) and poles of diagonal Padé
models ($\circ$) at $T=1$ (P1, P2) and $T=0.8$ (P3). The exact chains
continue to the left; the rational spectra are finite.}
\label{fig:polemaps}
\end{figure}

\begin{proposition}[A region containing all unstable roots]\label{prop:bound}
Let $p<m$, $D$ monic, $D(s)=s^m+\sum_{k<m}d_ks^k$, $N(s)=\sum_{k\le p}n_ks^k$,
and let $R>0$ be the positive root of
\[
  r^m-\sum_{k<m}|d_k|r^k-\sum_{k\le p}|n_k|r^k=0 .
\]
Then for every $T\ge0$, every root of~\eqref{eq:qp} with $\Real s\ge0$
satisfies $|s|\le R$.
\end{proposition}
The proof is in Appendix~\ref{ap:proofs}. For P1, $|s+1|\le2$ gives $R=3$;
for P3, $R\approx3.076$. The function \code{rhp\_root\_bound} computes $R$,
and \code{unstable\_root\_count} applies the argument principle to the
rectangle $[0,R+1]\times[-(R+1),R+1]$: the number $Z(T)$ of unstable roots
is obtained without guessing a search window and without approximation
(the count itself is a floating-point computation, reliable away from
critical delays).

\section{Roots for a given delay}

The study of this chapter locates roots with the delay-specific solver
\code{delay\_roots} (damped Newton from a grid, audited by an
argument-principle count); \code{croots} gives the same roots with the
subdivision algorithm of Chapter~\ref{ch:algorithm}. Table~\ref{tab:solver}
shows that the number of roots found agrees with the independent count and
that the residuals are near machine precision.

\input{../output/time_delay/data/table_solver.tex}

\section{Critical delays}

As $T$ varies the roots move continuously, and new roots appear only from
$\Real s=-\infty$. Stability can therefore change only when a root crosses
the imaginary axis. At a crossing $s=\ii\omega$, $\omega>0$,
\begin{equation}
  e^{-\ii\omega T}=-\frac{D(\ii\omega)}{N(\ii\omega)} .
  \label{eq:crossing}
\end{equation}
This complex equation splits into two real equations with very different
roles.

\paragraph{Modulus: where.} The left side has unit modulus for every $T$,
so the crossing frequencies do not depend on the delay and solve
\begin{equation}
  G(\omega)=|D(\ii\omega)|^2-|N(\ii\omega)|^2=0 ,
  \label{eq:G}
\end{equation}
an even polynomial in $\omega$, i.e.\ a polynomial of degree $m$ in
$x=\omega^2$. Its positive roots give a \emph{finite} set of frequencies.

\paragraph{Phase: when.} For each frequency $\omega_i$, with
$\theta_i=\operatorname{mod}(-\arg[-D(\ii\omega_i)/N(\ii\omega_i)],2\pi)$, the
phase condition gives \emph{infinitely many} critical delays
\begin{equation}
  T_{i,k}=\frac{\theta_i+2\pi k}{\omega_i},\qquad k=0,1,2,\ldots
  \label{eq:Tk}
\end{equation}
This elimination is the scalar form of the direct method
\citep{walton1987,olgac2002,michiels2014}. \code{critical\_delays} computes
all $T_{i,k}$ in an interval and refines each pair $(\omega,T)$ by a
two-variable Newton correction.

\begin{example}[P1]
$G(\omega)=1+\omega^2-4$ gives $\omega_c=\sqrt3$. Since
$-D(\ii\sqrt3)/N=e^{-2\pi\ii/3}$, $\theta=2\pi/3$ and
$T_k=(2\pi/3+2\pi k)/\sqrt3$: $T_0=2\pi/(3\sqrt3)\approx1.2092$,
$T_1\approx4.8368$, $T_2\approx8.4644$.
\end{example}
\begin{example}[P3]
$G(\omega)=\omega^4-7.36\,\omega^2+7$ has the roots $\omega_1\approx2.4976$ and
$\omega_2\approx1.0593$; each generates its own sequence of critical delays,
and the two sequences interleave.
\end{example}
For P2, $G(\omega)=\omega^2+3>0$: there is no crossing, and since P2 is
stable at $T=0$ it is stable for every delay.

\section{Direction of the crossings}

Differentiating $F(s(T),T)=0$ gives
\begin{equation}
  \frac{ds}{dT}=-\frac{F_T}{F_s}
  =\frac{s\,e^{-sT}N(s)}{D'(s)+e^{-sT}[N'(s)-TN(s)]} .
  \label{eq:speed}
\end{equation}
The crossing is destabilizing if $\Real(ds/dT)>0$ at $s=\ii\omega$. The
following classical result \citep{cooke1986} makes the direction a property
of the frequency alone.

\begin{proposition}[Direction depends only on the frequency]\label{prop:direction}
At a simple crossing $s=\ii\omega$, $\omega>0$,
$\operatorname{sign}\Real\dfrac{ds}{dT}=\operatorname{sign}G'(\omega)$.
\end{proposition}

All branches $k$ of a frequency therefore cross in the same direction. In
P1, $G'(\sqrt3)>0$: every $T_k$ adds an unstable pair. In P3, crossings at
$\omega_1$ always destabilize and crossings at $\omega_2$ always stabilize
(Table~\ref{tab:critical}, Figure~\ref{fig:critical}).

\input{../output/time_delay/data/table_critical.tex}

\begin{figure}[H]
\centering
\includegraphics[width=0.95\textwidth]{critical_delays.pdf}
\caption{Frequencies and directions of the exact crossings (up: into the
right half-plane; down: back to the left half-plane).}
\label{fig:critical}
\end{figure}

\section{Stability map}

Each simple crossing moves one conjugate pair, so the number of roots in
the right half-plane is the step function
\begin{equation}
  Z(T)=Z(0)+2\sum_{T_{i,k}<T}\operatorname{sign}G'(\omega_i),
  \label{eq:Z}
\end{equation}
where $Z(0)$ is the number of unstable roots of $D+N$ (the roots created
when $T$ becomes positive come from $\Real s=-\infty$). For P3 on
$[0,10]$:
\[
  Z(T)=\begin{cases}
   0, & 0\le T<0.2918,\\
   2, & 0.2918<T<2.6953,\\
   0, & 2.6953<T<2.8076\quad\text{(stability window)},\\
   2,4,6,4, & \text{after }2.8076,\ 5.3233,\ 7.8390,\ 8.6265.
  \end{cases}
\]
The window has width $0.112$ and could easily be missed by a sweep in $T$.
Formula~\eqref{eq:Z} was checked independently with
\code{unstable\_root\_count} at 80 delays, with full agreement (black curve
and circles in Figure~\ref{fig:Zp3}). Figure~\ref{fig:alpha} shows the
spectral abscissa computed by a sweep, refined near the window.

\begin{figure}[H]
\centering
\includegraphics[width=\textwidth]{spectral_abscissa.pdf}
\caption{Spectral abscissa versus delay, exact and with Padé models of
orders 1, 4 and 8; diamonds are the exact critical delays of~\eqref{eq:Tk}.
In P3 the exact curve touches zero at $T\approx2.695$ and crosses again at
$T\approx2.808$; the \pade{1} model misses the window and the \pade{4}
model widens it.}
\label{fig:alpha}
\end{figure}

\section{The Padé approximation}\label{sec:pade}

The diagonal Padé approximation of $e^{-z}$ of order $n$ is the rational
function $P_n(z)/Q_n(z)$ whose Taylor series agrees with that of $e^{-z}$
up to $z^{2n}$ \citep{baker1996}:
\begin{equation}
  e^{-z}-\frac{P_n(z)}{Q_n(z)}=O(z^{2n+1}),\quad
  Q_n(z)=\sum_{k=0}^n\frac{(2n-k)!\,n!}{(2n)!\,k!\,(n-k)!}z^k,\quad
  P_n(z)=Q_n(-z).
  \label{eq:pade}
\end{equation}
With $z=sT$ the characteristic equation becomes the polynomial
$D(s)Q_n(sT)+N(s)P_n(sT)$ of degree $m+n$. Three structural properties
determine its behavior.

\paragraph{(i) Exact gain.} Since $P_n(\ii y)=\overline{Q_n(\ii y)}$,
\begin{equation}
  \left|\frac{P_n(\ii y)}{Q_n(\ii y)}\right|=1\ \ (y\in\mathbb R),\qquad
  \left|\frac{P_n(z)}{Q_n(z)}\right|\le1\ \ (\Real z\ge0),
  \label{eq:allpass}
\end{equation}
the second by the maximum principle, as all roots of $Q_n$ lie in the left
half-plane. The error of the approximation is entirely in the phase.

\paragraph{(ii) Bounded phase.} Writing $P_n(\ii y)/Q_n(\ii y)=e^{\ii\phi_n(y)}$,
$\phi_n=-2\arg Q_n(\ii y)$ decreases monotonically and
\begin{equation}
  \phi_n(y)=-y+O(y^{2n+1})\ (y\to0),\qquad \phi_n(y)\searrow-n\pi\ (y\to\infty),
  \label{eq:phase}
\end{equation}
whereas the exact phase $-y$ is unbounded (Figure~\ref{fig:phaseuniv}).
This behavior depends only on $y=\omega T$.

\paragraph{(iii) Right-half-plane zeros.} The $n$ zeros of $P_n(sT)$ have
positive real part. The step response of the approximation jumps to
$(-1)^n$ at $t=0^+$ and oscillates before $t=T$, while the true delay stays
at zero (Figure~\ref{fig:step}).

\begin{figure}[H]
\centering
\includegraphics[width=\textwidth]{pade_phase_universal.pdf}
\caption{(a) Unwrapped phase of $e^{-\ii y}$ and of $\pade{n}$
approximations; dotted lines mark $-n\pi$. (b) Phase error; the dashed line
marks $10^{-2}$ rad.}
\label{fig:phaseuniv}
\end{figure}

\begin{figure}[H]
\centering
\includegraphics[width=0.85\textwidth]{pade_step_response.pdf}
\caption{Step response of $e^{-s}$ and of its Padé approximations.}
\label{fig:step}
\end{figure}

Table~\ref{tab:validity} quantifies property (ii): the largest $y$ for
which the phase error is below a tolerance. For $0.01$ rad and $n\ge4$,
$y_{0.01}\approx1.6n-2$.

\begin{table}[H]\centering\small
\caption{Largest $y=\omega T$ with Padé phase error below $0.1$, $0.01$ and
$10^{-6}$ rad, and asymptotic phase.}\label{tab:validity}
\input{../output/time_delay/data/pade_validity_table.tex}
\end{table}

\section{What Padé preserves}

\begin{proposition}[Frequencies and directions are inherited]\label{prop:padedir}
For the rational equation $D+NP_n/Q_n=0$:
\begin{enumerate}[label=(\alph*)]
  \item imaginary roots occur exactly at the frequencies $\omega_i$
        of~\eqref{eq:G};
  \item at $\omega_i$, the critical delays solve
        $\phi_n(\omega_iT)=-(\theta_i+2\pi k)$, which has a (unique)
        solution if and only if $\theta_i+2\pi k<n\pi$;
  \item every crossing has the direction $\operatorname{sign}G'(\omega_i)$.
\end{enumerate}
\end{proposition}
Diagonal Padé never invents crossing frequencies nor reverses a direction;
but a model of order $n$ has at most $\lceil(n\pi-\theta_i)/(2\pi)\rceil$
crossings per frequency for \emph{all} $T\ge0$, about $n/2$, where the true
system has infinitely many. Conclusions that depend only on the modulus are
preserved: for P2, $|s+2|\ge2>1\ge|P_n/Q_n|$ in the right half-plane, and every
Padé model is stable for every delay, as the true system.

In a fixed condition the dominant roots converge very fast with $n$
(Table~\ref{tab:pade}, Figure~\ref{fig:conv}): faster than geometrically,
consistent with the $O(y^{2n+1})$ error, until round-off. Far from the
origin, however, rational poles appear that have no exact counterpart (two
for $n=4$, four for $n=8$, six for $n=12$ in the window of
Figure~\ref{fig:polemaps}).

\input{../output/time_delay/data/table_pade_rms.tex}

\begin{figure}[H]
\centering
\includegraphics[width=0.7\textwidth]{pade_convergence.pdf}
\caption{RMS error of the dominant pair versus Padé order.}
\label{fig:conv}
\end{figure}

\section{What Padé loses}

\subsection{A non-conservative delay margin}

For P1, the \pade{1} phase is $-2\arctan(\omega T/2)$; equating it to
$-2\pi/3$ at $\omega_c=\sqrt3$ gives $T_c^{\pade{1}}=2$, against the exact
$1.2092$: a 65\,\% error (Table~\ref{tab:padeTc}). The rational polynomial
$\tfrac T2s^2+(1-\tfrac T2)s+3$ is Hurwitz if and only if $T<2$. The error
\emph{overestimates} the tolerable delay: at $T=1.5$, the \pade{1} model is
stable (abscissa $-1/6$) while the true system is unstable (abscissa
$+0.0656$). Figure~\ref{fig:sim} confirms it in the time domain, with the
delay equation integrated directly by \code{dde23} \citep{shampine2001}.
The first crossing occurs at $y=2\pi/3\approx2.09$, beyond the useful range
of orders 1 and 2 and within that of order 4 (Table~\ref{tab:validity}).

\begin{figure}[H]
\centering
\includegraphics[width=0.85\textwidth]{p1_time_simulation.pdf}
\caption{P1 at $T=1.5$, above the exact critical delay $1.209$: the delay
equation diverges; the \pade{1} model decays; \pade{4} coincides with the
exact solution.}
\label{fig:sim}
\end{figure}

\begin{remark}
The rational model also needs an initial state for the Padé filter. The
true initial condition is a function on $[-T,0]$; a model of order $n$ holds
only $n$ numbers. The constant history of Figure~\ref{fig:sim} can be
represented consistently; general histories cannot.
\end{remark}

\subsection{Shifted and missing branches}

For larger delays the branches $k\ge1$ of~\eqref{eq:Tk} matter, with larger
values of $y$. Figure~\ref{fig:padephase} shows the mechanism for P1: every
intersection of a phase curve with a horizontal line $-(\theta+2\pi k)$ is a
critical delay; the Padé phases saturate and reach only the first lines,
later and later. Table~\ref{tab:padeCoverage} and
Figures~\ref{fig:padeacc}--\ref{fig:padeevents} separate two notions of
quality: \emph{local accuracy} of the first critical delay converges fast
with $n$; \emph{global coverage} of the events grows slowly. In P1 there are
three events up to $T=10$; \pade{1} and \pade{2} represent one, \pade{4}
two and \pade{8} three. Even when an event is represented, its position may
be far off: \pade{4} places the second crossing of P1 at $5.71$ instead of
$4.84$.

\begin{figure}[H]
\centering
\includegraphics[width=0.85\textwidth]{pade_critical_phase.pdf}
\caption{P1: exact and Padé phases versus delay; intersections with the
branch lines are critical delays.}
\label{fig:padephase}
\end{figure}

\input{../output/time_delay/data/pade_critical_tables.tex}

\begin{figure}[H]
\centering
\includegraphics[width=\textwidth]{pade_critical_accuracy.pdf}
\caption{Left: relative error in the first critical delay. Right: fraction
of the exact crossings in $[0,10]$ represented by each order.}
\label{fig:padeacc}
\end{figure}

\begin{figure}[H]
\centering
\includegraphics[width=\textwidth]{pade_critical_events.pdf}
\caption{Critical delays predicted by each order (dots) and exact
(vertical lines and crosses).}
\label{fig:padeevents}
\end{figure}

\subsection{Stability windows erased, shifted and invented}

In P3 stability depends on the order of events at two frequencies, so a
shifted event changes the qualitative conclusion. Figure~\ref{fig:Zp3}
compares $Z(T)$ from~\eqref{eq:Z} with the rational models (their counts
computed with Proposition~\ref{prop:padedir} and checked by eigenvalues of
a balanced state-space realization), and Table~\ref{tab:intervals} lists
the predicted stable intervals.

\begin{figure}[H]
\centering
\includegraphics[width=\textwidth]{p3_unstable_count.pdf}
\caption{P3: number of right-half-plane roots. Black: exact~\eqref{eq:Z};
circles: independent bound-based count (Proposition~\ref{prop:bound}); red:
Padé model. Shaded: the rational model reaches the wrong stability
conclusion.}
\label{fig:Zp3}
\end{figure}

\begin{table}[H]\centering\small
\caption{P3: delay intervals in $[0,10]$ where each model predicts stability,
and the number of unstable roots at $T=10$.}\label{tab:intervals}
\input{../output/time_delay/data/p3_stability_intervals.tex}
\end{table}

Four kinds of error appear:
\begin{itemize}
  \item \textbf{Erased window.} \pade{1} has no stabilizing event of
        $\omega_2$ in $[0,10]$ and misses the window.
  \item \textbf{Unbounded fictitious stability.} \pade{2} has the
        stabilizing event at $T=2.887$ but not the second branch of
        $\omega_1$ (which needs $y\approx7.0>2\pi$); it predicts stability
        for \emph{every} $T>2.887$, whereas the system is unstable for every
        $T>2.808$.
  \item \textbf{Shifted and widened window.} \pade{3} and \pade{4} predict
        windows 12 and 2.5 times wider than the true one.
  \item \textbf{Invented stable region.} \pade{5} reproduces the window but,
        lacking the branches $k=2,3$ of $\omega_1$, turns the stabilizing
        event at $T=9.13$ into a stable region $(9.13,10]$ that does not exist.
\end{itemize}
Only from order 10 on is the stability map in $[0,10]$ correct, and any fixed
order fails again on a long enough interval. High orders are also
numerically fragile: the roots of the expanded polynomial
$DQ_n+NP_n$ gave a wrong unstable-root count for 24 of 99 small delays with
$n=10$, and 64 of 99 with $n=12$; a balanced realization avoids this.

\section{Guidelines}

\begin{keybox}{Using Padé safely}
\begin{enumerate}[leftmargin=1.5em]
  \item Decide stability on the original equation: frequencies from
        \eqref{eq:G}, delays from \eqref{eq:Tk}, the map \eqref{eq:Z}, and a
        bound-based count as a check.
  \item Use Padé as a local tool: design near an operating point, low-order
        simulation models, or Newton seeds.
  \item Choose the order from the largest relevant $\omega T$ with
        Table~\ref{tab:validity}; an order $n$ represents no phase beyond
        $n\pi$.
  \item Distrust stability windows and ``stable for all larger delays''
        predicted by rational models.
  \item Validate in the time domain near critical delays.
  \item Avoid high orders through expanded polynomials; for larger problems
        use methods designed for delay equations
        \citep{breda2005,breda2015,vyhlidal2009}.
\end{enumerate}
\end{keybox}

\section{Limitations}

The direct method above assumes a single constant delay, real
coefficients, simple crossings and a retarded system. Tangential or
multiple crossings require higher-order analysis \citep{cooke1986,michiels2014}.
For neutral systems ($p=m$) Proposition~\ref{prop:bound} fails, the root
chains approach vertical lines, and stability can be sensitive to
infinitesimal changes of the delay \citep{hale1993,michiels2014}. With
several delays, crossings form curves and surfaces
\citep{olgac2002,michiels2007}; matrix systems lead to nonlinear eigenvalue
problems \citep{breda2015,michiels2014}.
